\section{The finite installation of the congruence integer into PropGK}
\label{sec:prop-gk-installation-v046}

\status{This chapter proves the complete original implication
\(\mathrm{CNR3}\Longrightarrow(\mathrm{OneChar}\Longrightarrow
\mathrm{PropGK})\). RNC16 supplies its first premise. The original
single-character Selmer estimate \(\mathrm{OneChar}\) remains an
explicit premise; its rigidity and cohomological producers are not
proved here. The result discharges the finite implication task,
not the separate one-character producer.}

The source revision is
\texttt{7cbf7a0c32fca1badf05931879a37c3f8a2b03b4}.
We retain the original coefficient datum: for a rational prime
\(\ell\), \(O\subset\overline{\mathbb Q}_\ell\) is a finite
\(\mathbb Z_\ell\)-algebra which is a discrete valuation ring.
A Dirichlet character modulo \(N\ge1\) means a character of
\(G_N=(\mathbb Z/N\mathbb Z)^\times\), extended by zero off its
units. For \(H\le G_N\), set
\[
 \mathcal O_O(H)=\{\psi:G_N\to O^\times:\ker\psi=H\},\qquad
 \mathcal O_{\mathbb C}(H)=
       \{\chi:G_N\to\mathbb C^\times:\ker\chi=H\}.
 \tag{PGK.1}
\]
The first is exactly the original finite set \(\operatorname{orb}(c,H)\).

We also retain the actual source length, rather than substituting
an arbitrary nonnegative number. Let
\(\Psi=\varepsilon_\ell^e\psi_{\rm Gal}\), where
\(\psi_{\rm Gal}\) is the rational Galois character obtained by
acting on the \(N\)-th roots of unity. At an element of \(G_K\),
this character is restricted along the given map to \(G_{\mathbb Q}\).
Then
\(\mathcal L_\psi=V\!\operatorname{len}(\Omega,c,e,\psi)\) is the
\(O\)-length, possibly infinite, of the source Selmer module
\(\operatorname{ShaU}(\operatorname{twist}_2(\Psi^{-1}),\ell N,\top)\).
Its normalized real length is
\[
 v_\psi=V_n(\Omega,c,e,\psi)
       =\frac{\mathcal L_\psi.\operatorname{toNat}}{e_O},
 \qquad e_O=v_O(\ell)>0 .                                      \tag{PGK.2}
\]
We prove finiteness before using this as an ordinary length.
For a nonzero \(a\in\overline{\mathbb Q}_\ell\), write
\(v_\ell(a)=-\log\lVert a\rVert_\ell/\log\ell\), so
\(v_\ell(\ell)=1\).

Here is the exact retained \(\mathrm{OneChar}\) premise. For
every original \(\Omega,K,\ell,O,N,e,\psi\) satisfying
\(\mathrm{GD}\), every field isomorphism
\(\iota:\overline{\mathbb Q}_\ell\simeq\mathbb C\), every
\(\gamma\in G_K\), and
\[
 t=\Psi(\gamma)+1+\Psi(\gamma)^{-1}
                  \quad\hbox{in }\overline{\mathbb Q}_\ell,
 \tag{PGK.3}
\]
suppose \(\operatorname{tr}\rho(\gamma)\ne t\) for every actual
rank-three representation in
\(\mathrm{Adm3}(\iota,e,N)\) or
\(\mathrm{Adm2}(\iota,e,N|\operatorname{disc}K|)\).
Then \(\mathcal L_\psi<\infty\) and there is a finite set
\(X_\psi\) of such representations, pairwise nonisomorphic, with
\[
 v_\psi\le2\sum_{\rho\in X_\psi}
             v_\ell(\operatorname{tr}\rho(\gamma)-t).
 \tag{PGK.4}
\]
The scalar maps in (PGK.3) are the given inclusions
\(O\hookrightarrow\overline{\mathbb Q}_\ell\). The rank-two
alternative keeps the original two-by-two block in positions
\(0,2\) and its middle entry \(1\), as in RNC10. We assume exactly
this original premise, including its coefficient and local conditions.

% BEGIN PGK1
\paragraph{PGK1. The complete character orbit and its coefficient realization.}
\label{claim:PGK1}
Put \(h=[G_N:H]\). The complex set in (PGK.1) is empty or has
cardinality \(\varphi(h)\). If \(\psi_0\in\mathcal O_O(H)\) and
\(f:O\hookrightarrow\mathbb C\) is any injective ring map, then
\[
 E_f:\mathcal O_O(H)\overset\sim\longrightarrow
       \mathcal O_{\mathbb C}(H),\qquad
                     E_f(\psi)=f\circ\psi .                    \tag{PGK.5}
\]
Every member of \(\mathcal O_O(H)\) is \(\psi_0^j\) for some
\(j\) coprime to \(h\), and \(\operatorname{ord}(\psi_0)=h\).

\emph{Proof.}
A finite subgroup of the multiplicative group of a field is cyclic.
Indeed its exponent is attained by multiplying elements of maximum
prime-power order, and the polynomial \(X^d-1\) has at most \(d\)
roots; an exponent smaller than the group's order is impossible.
A character with kernel \(H\) therefore identifies \(G_N/H\)
with a cyclic group of order \(h\). Fix a generator \(\bar g\)
and put \(\zeta=\psi_0(g)\). Its order is \(h\).
The \(h\) powers of \(f(\zeta)\) are all roots of \(X^h-1\)
in \(\mathbb C\), by degree and distinctness. Any faithful
character of \(G_N/H\) consequently sends \(\bar g\) to
\(f(\zeta)^j\) with \(\gcd(j,h)=1\). It is precisely
\(f\circ\psi_0^j\). Such powers already take their values in
the given \(O\); no enlargement of the coefficient ring is needed.
They have distinct values for distinct invertible residue classes
\(j\pmod h\), and all have kernel \(H\).
Injectivity of \(f\) also proves that composition preserves the
kernel and distinguishes characters. This proves (PGK.5), its
cardinality, and the assertion about powers. If the complex set
has no member its cardinality is zero. The argument includes
\(h=1\), where there is one trivial character and
\(\varphi(1)=1\). \hfill\(\square\)
% END PGK1

% BEGIN PGK2
\paragraph{PGK2. Every character with the same kernel has the original GD data.}
\label{claim:PGK2}
If \(\mathrm{GD}(\Omega,K,\ell,e,\psi_0)\) and
\(\ker\psi=\ker\psi_0\), then
\(\mathrm{GD}(\Omega,K,\ell,e,\psi)\).

\emph{Proof.}
We check all the fields of the original record. Its conditions are
\[
 \begin{gathered}
 \ell\ne2,\quad e\ge1,\quad 2\mid N,\quad
 \ell\nmid\operatorname{ord}\psi_0,\quad
       (\operatorname{comp}_\ell\psi_0)^2=1,\quad
       \psi_0(-1)=(-1)^e,\\
 [K:\mathbb Q]=2,\quad K\text{ totally complex},\quad
 \ell\nmid|\operatorname{disc}K|\,h_K,\quad
       (|\operatorname{disc}K|,N)=1,\quad \ell-1\nmid2e,\\
 \ker\bar\varepsilon_\ell\cap
     \ker(\operatorname{chiP}_\ell(\psi_0)_{\rm Gal})
                                      \not\subseteq G_K .
 \end{gathered}                                                  \tag{PGK.6}
\]
Here \(h_K\) is the class number, \(G_K\) denotes the subgroup
fixing the given copy of \(K\) in \(\Omega\), and
\(\operatorname{comp}_\ell\) is restriction to the
\(\ell\)-primary factor of the modulus, followed by inflation.
Equivalently it is composition with the original Chinese-remainder
idempotent on \(G_N\), so it commutes with products and powers.

By PGK1 write \(\psi=\psi_0^j\), \(\gcd(j,h)=1\), where
\(h=\operatorname{ord}\psi_0\). Its order remains \(h\).
Since \(\ell\nmid h\), the source's \(\ell\)-power part
\(\operatorname{kap}_\ell(\psi_0)\) is trivial and its
prime-to-\(\ell\) part \(\operatorname{th0}_\ell(\psi_0)\)
is \(\psi_0\); the same holds for \(\psi\). This also follows
directly from their defining exponents: the \(\ell\)-part of
\(h\) is \(1\), so the defining power is a multiple of \(h\).

Set \(\delta=\operatorname{comp}_\ell\psi_0\). Both
\(\delta^2=1\) and \(\delta^h=1\). Any element of order dividing
both \(2\) and \(h\) is fixed by raising to \(j\): if its order
is \(2\), then \(h\) is even and \(j\) odd. Thus
\(\operatorname{comp}_\ell\psi=\delta^j=\delta\), preserving
the square condition. The same argument applied to
\(\psi_0(-1)\) preserves the parity condition.
Furthermore
\[
 \operatorname{chiP}_\ell(\psi_0)=\psi_0\delta^{-1},\qquad
 \operatorname{chiP}_\ell(\psi)
                 =(\psi_0\delta^{-1})^j .                      \tag{PGK.7}
\]
The order of \(\psi_0\delta^{-1}\) divides \(h\), so taking
its \(j\)-th power preserves its kernel, also after composition
with the rational cyclotomic action. The last condition in
(PGK.6) is therefore unchanged. Every remaining displayed
condition depends only on \(K,\ell,e,N\) or on the order
already checked. This proves the whole record.
No additional prime bound has been imposed. In particular
\(\ell=3\) is included in the later universal quantifier;
its original \(\ell-1\nmid2e\) antecedent is impossible.
\hfill\(\square\)
% END PGK2

% BEGIN PGK3
\paragraph{PGK3. The precise Frobenius and Eisenstein scalar.}
\label{claim:PGK3}
Let \(p\nmid N\), \(\ell\ne p\), \(w\mid p\) in \(K\), and
\(q=\operatorname{Norm}(w)\). Choose a prolongation \(B\) to
\(\Omega\) and an arithmetic Frobenius \(\sigma\in G_K\) at
\(B\). For the image \(g\in G_{\mathbb Q}\),
\[
 \operatorname{cls}_N(g)=q\pmod N,\qquad
           \varepsilon_\ell(g)=q\quad\hbox{in }\mathbb Z_\ell^\times.
 \tag{PGK.8}
\]
For any \(\psi\), any \(\iota\), and
\(\chi=\iota\circ\psi\), put \(z_\chi=\chi(q)^{-1}\).
Then the original scalar satisfies
\[
 \begin{split}
 \iota^{-1}\bigl(q^{-e}z_\chi+1+q^ez_\chi^{-1}\bigr)
 &=\Psi(g^{-1})+1+\Psi(g^{-1})^{-1}\\
 &=q^{-e}\psi(q)^{-1}+1+q^e\psi(q).
 \end{split}                                                     \tag{PGK.9}
\]

\emph{Proof.}
The residue field at \(w\) has cardinality \(q=p^f\) for
some \(f\ge1\), hence \((q,N)=1\). For every positive integer
\(m\) prime to \(p\), roots of unity of order dividing \(m\)
are unramified at \(p\), and reduction is injective on them.
One way to see both statements is that \(X^m-1\) is separable
on the residue field and all its roots lift uniquely over the
maximal unramified extension. Arithmetic Frobenius raises
their reductions, and therefore the roots themselves, to
the \(q\)-th power. Apply this first to \(m=N\), then to
every \(m=\ell^a\), and take the inverse limit. This proves
(PGK.8) with the original arithmetic convention, including
its restriction from \(K\) to \(\mathbb Q\).

Consequently \(\Psi(g)=q^e\psi(q)\). Evaluating the
inverse character at \(g^{-1}\) gives (PGK.9), since a field
isomorphism fixes rational numbers. All quantities on the
right use the original inclusion of \(O\), without a change
of coefficient datum.

The choices used here exist for every later datum. Valuations
extend to an algebraic closure, and the local decomposition
group surjects onto the absolute Galois group of the finite
residue field, giving a lift of its arithmetic Frobenius.
Also \(\Omega\), already an algebraic closure of \(\mathbb Q\)
containing \(K\), is an algebraic closure of \(K\). A field
isomorphism \(\overline{\mathbb Q}_\ell\simeq\mathbb C\)
exists: both algebraically closed characteristic-zero fields
have cardinality and transcendence degree over \(\mathbb Q\)
equal to the continuum; match transcendence bases and extend
to their algebraic closures. No continuity of this isomorphism
is asserted or required. \hfill\(\square\)
% END PGK3

% BEGIN PGK4
\paragraph{PGK4. A single integer before all coefficient choices.}
\label{claim:PGK4}
Assume \(K\) is imaginary quadratic. There are natural numbers
\(c,A\), depending only on \(K\), such that for all
\(e,N\ge1\), \(H\le G_N\), and prime \(p\) with
\(p\nmid30N|\operatorname{disc}K|\), there is
\(NN\in\mathbb Z\setminus\{0\}\) with
\[
 \log|NN|\le
       c\varphi([G_N:H])\,[N(e+1)]^A\log p.                    \tag{PGK.10}
\]
It simultaneously has the trace-separation and norm-product
properties of RNC16 at every later original coefficient datum,
with the index set \(\mathcal O_{\mathbb C}(H)\) and values
\(z_\chi=\chi(q)^{-1}\).

\emph{Proof.}
Choose \(c,A\) from RNC16 for this \(K\). The prime restriction
implies \(p\ge7\), \(p\nmid N\), and
\(p\nmid\operatorname{disc}K\): the smaller primes are
\(2,3,5\), all dividing \(30\). Choose any place \(w\mid p\)
and write \(q=\operatorname{Norm}(w)\).
The finite set \(\mathcal O_{\mathbb C}(H)\) is determined
before \(\ell\) and \(O\). Each \(\chi(q)\) is a root of
unity, because \((q,N)=1\). An automorphism \(a\) of
\(\mathbb C\) sends \(\chi\) to \(a\circ\chi\), preserving
its kernel and permuting this finite set, with inverse given
by \(a^{-1}\). It sends \(z_\chi\) to \(z_{a\circ\chi}\).
Thus this is exactly an admissible indexed orbit for RNC16,
including repetitions of the value at \(q\).

Apply RNC16 to obtain \(NN\ne0\), its full product inequality,
and the logarithmic bound with
\(|\mathcal O_{\mathbb C}(H)|\) in place of \(\varphi([G_N:H])\).
PGK1 bounds that cardinality by the latter number. All remaining
factors in the bound are nonnegative, giving (PGK.10).
This invocation precedes every choice of \(\ell,\Omega,O,\iota,B\),
Frobenius, or finite collection of representations. An empty
complex orbit causes no difficulty; the original RNC16 product
is then the empty product \(1\). \hfill\(\square\)
% END PGK4

% BEGIN PGK5
\paragraph{PGK5. Installation of the original single-character estimates.}
\label{claim:PGK5}
Retain the integer from PGK4 and assume the exact
\(\mathrm{OneChar}\) premise (PGK.3)--(PGK.4). For every later
\(\Omega,\ell\ne p,O\) such that at least one
\(\psi_0\in\mathcal O_O(H)\) has \(\mathrm{GD}\), all the
original lengths \(\mathcal L_\psi\),
\(\psi\in\mathcal O_O(H)\), are finite. There are finite
pairwise-nonisomorphic sets \(X_\psi\) of the original admissible
representations for which, writing
\[
 t_\psi=\Psi(g^{-1})+1+\Psi(g^{-1})^{-1},\qquad
 a_{\psi,\rho}=\operatorname{tr}\rho(\sigma^{-1})-t_\psi,
 \tag{PGK.11}
\]
one has
\[
 a_{\psi,\rho}\ne0,\qquad
 v_\psi\le2\sum_{\rho\in X_\psi}v_\ell(a_{\psi,\rho}),\qquad
 \lVert NN\rVert_\ell\le
       \prod_{\psi\in\mathcal O_O(H)}\prod_{\rho\in X_\psi}
                         \lVert a_{\psi,\rho}\rVert_\ell .      \tag{PGK.12}
\]

\emph{Proof.}
The assumed \(\mathrm{GD}\) gives \(\ell\ne2\), so RNC16
applies to this coefficient prime, without an additional
restriction. Choose \(\iota,B,\sigma\) as in PGK3. Use PGK1
with the injective map
\(O\hookrightarrow\overline{\mathbb Q}_\ell
\xrightarrow{\iota}\mathbb C\) to obtain the bijection \(E\)
of the two complete character orbits. PGK2 supplies
\(\mathrm{GD}\) for every \(\psi\) in the coefficient orbit.

For each such \(\psi\), (PGK.9) identifies \(t_\psi\) with
\(\iota^{-1}E(q,e,z_{E\psi})\), using exactly the Eisenstein
expression in RNC16. Its universal nonvanishing conclusion
gives the antecedent in (PGK.3) for \(\gamma=\sigma^{-1}\).
To see that this precedes the choice of a finite set, apply
RNC16 with the singleton \(\{\rho\}\) at each orbit index
for an arbitrary admissible \(\rho\). Each such family meets
its pairwise-nonisomorphism condition. Thus every actual
admissible \(\rho\) has trace different from this scalar.

The stated \(\mathrm{OneChar}\) premise now yields
\(\mathcal L_\psi<\infty\) and a set \(X_\psi\) satisfying
its exact admissibility, nonisomorphism and estimate. These
choices are made for finitely many \(\psi\). Apply the product
clause of the same RNC16 integer to the family
\(X_{E^{-1}\chi}\) indexed by
\(\chi\in\mathcal O_{\mathbb C}(H)\). Its hypotheses are
exactly the properties just supplied by \(\mathrm{OneChar}\).
Reindexing through \(E\) and using (PGK.9) gives the last
inequality of (PGK.12), and its nonvanishing clause gives the
first. A representation may occur for several characters;
the original product retains its character index, so no
cross-character nonisomorphism is required. \hfill\(\square\)
% END PGK5

% BEGIN PGK6
\paragraph{PGK6. The whole original PropGK implication.}
\label{claim:PGK6}
Assuming the original \(\mathrm{OneChar}\), RNC16 proves the
literal \(\mathrm{PropGK}\): for every number field \(K\) there
are \(c_5,A\in\mathbb N\) such that, for every
\(e,N\ge1\), every \(H\le G_N\), and every prime
\(p\nmid30N|\operatorname{disc}K|\), a nonzero integer \(NN\)
satisfies (PGK.10) with \(c=c_5\), and, for every later
\(\Omega,\ell\ne p,O\),
\[
 \begin{split}
 &(\exists\psi_0\in\mathcal O_O(H),\ \mathrm{GD}(\psi_0))\\
 &\quad\Longrightarrow
   \left(\mathcal L_\psi<\infty\text{ for every }\psi\in
                        \mathcal O_O(H),\quad
     \sum_{\psi\in\mathcal O_O(H)}v_\psi
                           \le2v_\ell(NN)\right).
 \end{split}                                                     \tag{PGK.13}
\]
Here \(v_\ell(NN)=\operatorname{padicValInt}(\ell,NN)\).

\emph{Proof.}
First suppose \(K\) is imaginary quadratic. Choose the
constants and the integer by PGK4. For any later coefficient
datum with the displayed antecedent, PGK5 gives all the
finiteness assertions and (PGK.12). Every factor there is
strictly positive in norm, and \(NN\ne0\). We may therefore
take logarithms to base \(\ell>1\), obtaining
\[
 -v_\ell(NN)\le
       -\sum_{\psi\in\mathcal O_O(H)}\sum_{\rho\in X_\psi}
                            v_\ell(a_{\psi,\rho}).
 \tag{PGK.14}
\]
The identity on the left follows directly from the usual
integer factorization \(NN=\ell^r u\), \(\ell\nmid u\):
its norm is \(\ell^{-r}\). The logarithm of each finite
product is the sum of the logarithms, also for empty sets.
Multiplying (PGK.14) by \(-1\), summing (PGK.4), and
multiplying by \(2\) gives
\[
 \sum_\psi v_\psi
 \le2\sum_\psi\sum_{\rho\in X_\psi}v_\ell(a_{\psi,\rho})
 \le2v_\ell(NN).
\]
This proves exactly (PGK.13) for the same integer chosen
before the coefficient prime. No sign assumption on the
individual valuations was needed for the logarithmic step.

For any number field \(K\) which is not totally complex of
degree two, choose \(c_5=A=0\) and \(NN=1\) for all the
varying data. The logarithmic bound is \(0\le0\). The later
antecedent is impossible, since the \(\mathrm{quad}\) and
\(\mathrm{imag}\) fields in (PGK.6) assert precisely the
excluded conjunction. A totally complex quadratic field is
CM over \(\mathbb Q\), so these two cases exhaust all
number fields. This completes the literal original quantifier
order and the whole finite implication. Its only unsupplied
arithmetic premise is the original \(\mathrm{OneChar}\);
in particular this proof does not establish its type-II
rigidity producer. \hfill\(\square\)
% END PGK6

\subsection{Scope of the completed implication}

The construction has order
\[
 K\ \longmapsto\ (c_5,A)\ \longmapsto\ (e,N,H,p)
 \ \longmapsto\ NN\ne0\ \longmapsto\ (\Omega,\ell,O).
\]
The coefficient orbit and the sets in the one-character
estimate are chosen only after \(NN\). The two complete
orbits are identified by (PGK.5), and the arithmetic versus
geometric Frobenius sign is fixed by (PGK.9). Thus the
finite orbit/Frobenius/valuation task is proved with its
original hypotheses. The separate production of
\(\mathrm{OneChar}\), including the original
\(\mathrm{RigII}\) and other arithmetic inputs, remains outside
this component.
